\section*{Appendix: Details of SENet-154} \label{app:senet154}

SENet-154 is constructed by incorporating SE blocks into a modified version of the 64$\times$4d ResNeXt-152 which extends the original ResNeXt-101 \cite{xie_cvpr2017resnext} by adopting the block stacking strategy of ResNet-152 \cite{he_cvpr2016resnet}. Further differences to the design and training of this model (beyond the use of SE blocks) are as follows:
(a) The number of the first $1\times1$ convolutional channels for each bottleneck building block was halved to reduce the computational cost of the model with a minimal decrease in performance.
(b) The first $7\times7$ convolutional layer was replaced with three consecutive $3\times3$ convolutional layers.
(c) The $1\times1$ down-sampling projection with stride-$2$ convolution was replaced with a $3\times3$ stride-$2$ convolution to preserve information.
(d) A dropout layer (with a dropout ratio of $0.2$) was inserted before the classification layer to reduce overfitting.
(e) Label-smoothing regularisation (as introduced in \cite{szegedy_cvpr2016inceptionv3}) was used during training.
(f) The parameters of all BN layers were frozen for the last few training epochs to ensure consistency between training and testing.
(g) Training was performed with 8 servers (64 GPUs) in parallel to enable large batch sizes (2048).  The initial learning rate was set to $1.0$.

\begin{figure}[t]
\begin{center}
\subfigure[goldfish]{\includegraphics[height=16mm]{figures/goldfish.JPEG}}  \hspace{0.5mm}
\subfigure[pug]{\includegraphics[height=16mm]{figures/pug.JPEG}} 
    \hspace{0.5mm}
\subfigure[plane]{\includegraphics[height=16mm]{figures/plane.JPEG}} 
    \hspace{0.5mm}
\subfigure[cliff]{\includegraphics[height=16mm]{figures/cliff.JPEG}}
    \hspace{0.5mm}
\end{center}
\vspace{-1.5em}
\caption{Sample images from the four classes of ImageNet used in the experiments described in Sec.~7.2.%~\ref{subsection:role-of-excitation}.
}
\label{fig:sample-ims}
\end{figure}



